\subsection{紧度量空间；紧可度量化空间}

一致空间预紧性的判别法（第二章，§ 4，第 2 号，定理 3）对度量空间给出下述命题：

命题 14. 度量空间 X 预紧的必要充分条件是：对每个 \(\varepsilon > 0\) ，存在 X 的由直径 \(\leqslant \varepsilon\) 的集组成的有限覆盖。

若再加上 X 完备这一假设，我们就得到度量空间紧性的一个判别法。

由命题 14 我们得到一个适用于可度量化空间的紧性拓扑判别法：

命题 15. 可度量化拓扑空间 X 紧的必要充分条件是：X 的每个无穷点列在 X 中有一个聚点。

第一章，§ 9，第 1 号的公理 (C) 表明该条件是必要的。为证明充分性，设 \(d\) 是与 \(\mathbf{X}\) 的拓扑相容的一个度量。我们先证明这样定义的度量空间 \(\mathbf{X}\) 是完备的：\(\mathbf{X}\) 中每个柯西序列都有聚点，从而是收敛的（第二章，§ 3，第 2 号，命题 5 的推论 2）；因而由命题 9，\(\mathbf{X}\) 是完备的。其次我们证明 \(\mathbf{X}\) 是预紧的；若非如此，则由命题 14，存在实数 \(\alpha > 0\) ，使得 \(\mathbf{X}\) 不能被 \(\mathbf{X}\) 的任何有限多个直径 \(\leqslant \alpha\) 的子集覆盖。于是我们可以对 \(n\) 用归纳法定义一个无穷序列 \((x_n)\) （\(\mathbf{X}\) 的点的序列），满足条件 \(d(x_p, x_n) > \frac{1}{2} \alpha\) 对一切 \(p < n\) 成立；而这样的序列不可能有聚点，因为每个半径 \(< \frac{1}{2} \alpha\) 的球至多包含该序列的一个点。

推论. 可度量化拓扑空间 X 的子集 A 相对紧的必要充分条件是：A 的每个无穷点列在 X 中有一个聚点。

设 \(d\) 是与 \(\mathbf{X}\) 的拓扑相容的一个度量。我们将应用命题 15 的判别法证明空间 \(\overline{\mathbf{A}}\) 是紧的。设 \((x_{n})\) 是 \(\overline{\mathbf{A}}\) 的点的一个序列；则对每个指标 \(n\) 存在 \(y_{n} \in \mathbf{A}\) 使得 \(d(x_{n}, y_{n}) < 1/n\) ；由假设，序列 \((y_{n})\) 有一个聚点 \(a \in \mathbf{X}\) ，且 \(a\) 也是序列 \((x_{n})\) 的聚点，因为若 \(y_{m}\) 位于以 \(a\) 为中心、以 \(1/n\) 为半径的球中（对某个 \(m > n\) ），则 \(x_{m}\) 位于以 \(a\) 为中心、以 \(2/n\) 为半径的球中。

应当指出，命题 15 并不是 X 的每个点处可数基本邻域系存在性的推论；存在不可度量化的、非紧的空间的例子，其中每个点都有可数的基本邻域系，且每个无穷点列都有聚点（习题 15）。

命题 16. 紧空间 X 可度量化的必要充分条件是它的拓扑具有可数基。

该条件是必要的。由命题 14，对每个整数 \(n \geqslant 1\) ，存在一个有限子集 \(A_n\) （\(X\) 的子集），使得 \(A_n\) 与 \(X\) 的每个点的距离 \(\leqslant 1/n\) ；因而可数集 \(A = \bigcup_{n} A_n\) 在 \(X\) 中稠密，于是由第 8 号的命题 12 即得结论。

该条件是充分的。设 \((\mathbf{U}_{n})\) 是 X 的拓扑的一个可数基。因而对角线 \(\Delta\) —— \(X \times X\) 的对角线 —— 的每个点的每个邻域都包含一个形如 \(U_{n} \times U_{n}\) 的邻域；对紧子集 \(\Delta\) （ \(X \times X\) 的紧子集）应用波莱尔–勒贝格公理可知，\(\Delta\) 的每个邻域都包含形如 \(U_{n} \times U_{n}\) 的集的一个有限并，而这样的并是 \(\Delta\) 的一个邻域。因而 \(\Delta\) 的这样的邻域——由形如 \(U_{n} \times U_{n}\) 的集的有限并组成的邻域——构成 X 的一致结构的一个围域基本系（第二章，§ 4，第 1 号，定理 1），于是由第 4 号的定理 1 即得结论。

推论. 设 X 是局部紧空间，X' 是通过给 X 添加一个无穷远点 ω 而得到的紧空间（第一章，§ 9，第 8 号）。则下列陈述等价：

\begin{enumerate}
\def\labelenumi{\alph{enumi})}
\item
  \(\mathbf{X}\) 的拓扑有可数基。
\item
  \(\mathbf{X}'\) 可度量化。
\item
  X 可度量化且是 \(\sigma\) -紧的。
\end{enumerate}

a) \(\Longrightarrow\) b)：设 \((\mathbf{U}_n)\) 是 \(\mathbf{X}\) 的拓扑的一个可数基。点 \(x \in \mathbf{X}\) 的每个邻域都包含 \(x\) 的一个紧邻域，而后者又包含 \(x\) 的一个等于某个 \(\mathbf{U}_n\) 的邻域。因而相对紧的诸 \(\mathbf{U}_n\) 构成 \(\mathbf{X}\) 的拓扑的一个基，于是我们可以假定所有 \(\mathbf{U}_n\) 都是相对紧的。因而 \(\mathbf{X}\) 是紧集 \(\overline{\mathbf{U}}_n\) 的可数并，即它是 \(\sigma\) -紧的；这蕴含在 \(\mathbf{X}'\) 中点 \(\omega\) 有由开邻域组成的可数基本系 \((\mathbf{V}_n)\) （第一章，§ 9，第 9 号，命题 15 的推论 2）。因而点 \(y \in \mathbf{X}'\) 的每个邻域都包含某个 \(\mathbf{U}_n\) 或某个 \(\mathbf{V}_n\) ，它是 \(y\) 的一个邻域，于是诸 \(U_{n}\) 与诸 \(V_{n}\) 构成 \(X'\) 的拓扑的一个可数基。因而由命题 16，\(X'\) 可度量化。

b) \(\Longrightarrow\) c): 若 \(X'\) 可度量化，则 \(\omega\) 有可数的基本邻域系，因而由第一章，§ 9，第 9 号，命题 15 的推论 2，\(X\) 是 \(\sigma\) -紧的。

c) \(\Longrightarrow a)\) ：由假设，存在相对紧开集的一个递增序列 \((\mathbf{V}_n)\) ，它们覆盖 \(\mathbf{X}\) 且满足 \(\overline{\mathbf{V}}_n \subset \mathbf{V}_{n+1}\) （第一章，§ 9，第 9 号，命题 15）。子空间 \(\overline{\mathbf{V}}_n\) 紧且可度量化，因而具有可数基（命题 16），从而 \(\mathbf{V}_n\) 也具有可数基。设 \((\mathbf{U}_{mn})_{m \geqslant 1}\) 是 \(\mathbf{V}_n\) 的拓扑的一个基。对每个 \(x \in \mathbf{X}\) 与每个邻域 \(W\) （ \(x\) 的），存在 \(n\) 使得 \(x \in \mathbf{V}_n\) ，因而存在 \(m\) 使得 \(x \in \mathbf{U}_{mn} \subset \mathbf{V}_n \cap W\) 。因而诸集 \(\mathbf{U}_{mn}\) （ \(m \geqslant 1\) ， \(n \geqslant 1\) ）构成 \(\mathbf{X}\) 的拓扑的一个基。
% semantic-fragments: legacy-input-seam
\relax{}
